% TOP_PROBLEM: {"id": 3043, "title": "3-accessibility of Fibonacci numbers", "category_id": 2, "category": {"id": 2, "name": "combinatorics", "display_name": "Combinatorics", "description": "Counting problems, graph theory, discrete structures.", "slug": "combinatorics", "order_index": 2, "created_at": "2026-07-31T15:26:25.671Z"}, "status": "open", "problem_number": "OPG-1825", "set_id": 12}
% FIRST_POSTED: 2026-10-01
% ENTRY_KIND: statement_only
% UnsolvedMath ID 3043; source catalogue code OPG-1825
% !TeX program = lualatex
% Definition and sources only; this file is not an LLM solution attempt.
\documentclass[11pt,a4paper]{article}
\usepackage[margin=25mm]{geometry}
\usepackage{fontspec}
\setmainfont{lmroman10-regular.otf}[BoldFont=lmroman10-bold.otf,
 ItalicFont=lmroman10-italic.otf,BoldItalicFont=lmroman10-bolditalic.otf]
\usepackage{amsmath,amssymb,mathtools}
\usepackage{unicode-math}
\setmathfont{latinmodern-math.otf}
\AtBeginDocument{\let\setminus\smallsetminus}
\usepackage{xurl,hyperref}
\hypersetup{colorlinks=true,urlcolor=blue}
\setcounter{secnumdepth}{0}
\setlength{\parindent}{0pt}
\setlength{\parskip}{5pt}
\setlength{\emergencystretch}{3em}
\begin{document}
\section*{3-accessibility of Fibonacci numbers}\label{3043}
{\small UnsolvedMath ID 3043 · OPG-1825 · Combinatorics · OpenGarden}
\subsection{Definitions and mathematical statement}
Define $F_1=F_2=1$ and $F_{j+2}=F_{j+1}+F_j$ for $j\ge1$, and put
$\mathcal F=\{F_j:j\ge1\}\subseteq\mathbb N$. This is a set, so the two
initial occurrences of $1$ have no multiplicity significance. In particular,
$0$ is not an allowed gap.

Given a set $S\subseteq\mathbb N$, an $S$-difference sequence of length
$k$ is a strictly increasing sequence $x_1<\cdots<x_k$ of positive integers
such that $x_{i+1}-x_i\in S$ for every $1\le i<k$.
For a colouring $\chi:\mathbb N\to[r]$, the sequence is monochromatic if
$\chi(x_1)=\cdots=\chi(x_k)$. The set $S$ is called $r$-accessible when
for every such colouring and every integer $k\ge2$, a monochromatic
$S$-difference sequence of length $k$ exists.

Is $\mathcal F$ three-accessible? Written with its quantifiers, the target is
\[
 \forall\chi:\mathbb N\to\{1,2,3\}\quad\forall k\ge2\quad
 \exists x_1<\cdots<x_k:
 \ \chi(x_i)=\chi(x_1),\ \ x_{i+1}-x_i\in\mathcal F.
\]
The gap may vary from one step to the next, and Fibonacci values may be
reused as gaps. The terms $x_i$ themselves need not be Fibonacci numbers.
The length is arbitrary but finite; the statement does not require one
infinite monochromatic sequence. A negative answer must provide a
three-colouring and a fixed finite length which no monochromatic
Fibonacci-difference sequence reaches.
\subsection{Short English statement}
Does every three-colouring of the positive integers contain arbitrarily long monochromatic increasing sequences with Fibonacci gaps?
\subsection{Sources}
\begin{itemize}
\item[S1] ULAM.AI and the UnsolvedMath Contributors, supplied problems.json, record 3043 (OPG-1825), OpenGarden collection. Catalogue identity and source transcription; CC BY 4.0.
\url{https://huggingface.co/datasets/ulamai/UnsolvedMath}.
\item[S2] Open Problem Garden, 3-accessibility of Fibonacci numbers. Original problem and discussion; the mathematical formulation below is expanded from the supplied source record.
\url{http://www.openproblemgarden.org/op/3_accessibility_of_fibonacci_numbers}.
\end{itemize}
\end{document}
